% Generated draft: requires mathematical and source review.
\documentclass[11pt]{article}
\usepackage[margin=25mm]{geometry}
\usepackage{fontspec}
\setmainfont{Latin Modern Roman}
\usepackage{amsmath,amssymb,mathtools,unicode-math}
\setmathfont{Latin Modern Math}
\usepackage{enumitem,needspace,xurl,hyperref}
\setlength{\parindent}{0pt}
\setlength{\parskip}{5pt}
\setlength{\emergencystretch}{3em}
\newcommand{\N}{\mathbb{N}}
\newcommand{\Z}{\mathbb{Z}}
\newcommand{\Q}{\mathbb{Q}}
\newcommand{\R}{\mathbb{R}}
\newcommand{\C}{\mathbb{C}}
\newcommand{\F}{\mathbb{F}}
\newcommand{\A}{\mathbb{A}}
\newcommand{\PP}{\mathbb{P}}
\newcommand{\E}{\mathbb{E}}
\newcommand{\eps}{\varepsilon}
\newcommand{\dd}{\,\mathrm{d}}
\begin{document}
\section*{UnsolvedMath ID: 30000203}
\textbf{Source problem code:} catalogue entry
\subsection{Definitions and mathematical statement}
Let $G$ be a finite primitive permutation group in the O'Nan--Scott twisted-wreath (TW) class, with the particular decomposition from the source's Theorem 3
\[
 G=\operatorname{Soc}(G)\rtimes P\ \leq\ H\wr S_k
\]
in product action on $\Delta^k$, where $H\leq\operatorname{Sym}(\Delta)$ is the specified primitive component. The socle is the subgroup generated by all minimal normal subgroups; in the TW situation it is a regular nonabelian direct power of a simple group, and a point stabilizer is a complement represented here by $P$. Regular means transitive with trivial stabilizers in the socle action. Primitivity means transitivity with no nontrivial invariant partition.

For a finite transitive group $L$ on $\Lambda$, fix $\alpha\in\Lambda$ and define the minimal subdegree, in the orbit-away-from-the-basepoint convention, by
\[
 \operatorname{MinSubDeg}(L)=
 \min_{\beta\in\Lambda\setminus\{\alpha\}}|\beta^{L_\alpha}|.
\]
It is independent of the basepoint. For the nonregular primitive actions at issue, other point stabilizer orbits are nontrivial; if a different convention excludes all size-one orbits, that convention must be matched to the missing source definitions.

The question, called Question 4 in the supplied assessment, is whether the bound of Theorem 1 is always strict in this TW setting:
\[
 \operatorname{MinSubDeg}(G)
       < k\,\operatorname{MinSubDeg}(H).
\]
Thus the target excludes equality with the stated upper bound, not merely violation of a non-strict estimate. The integer $k$, component $H$, complement $P$ and wreath embedding must be those supplied by Theorem 3. The theorem itself, its restrictions on $k$ and its construction of these data are absent. The quantifier is not silently broadened to arbitrary wreath embeddings or arbitrary choices of $H$; in particular the corrected standalone sentence does not recover all qualifications in the original ``$P,H$ as in Theorem 3''.

The supplied literature triage checked 21 August 2026 reports bounds, constructions of minimal suborbits and varied twisted-wreath behavior in the primary study, but says the universal strict inequality remains a question there and that no complete resolution was found. These are dated supplied findings, not an independent current-status assessment.
\subsection{Short English statement}
For the source's specified product-action realization of a primitive twisted-wreath group, must its smallest nontrivial suborbit be strictly smaller than $k$ times that of its primitive component? The missing Theorem 3 is needed to specify the permitted decomposition; the supplied assessment leaves the universal strict bound unresolved.
\subsection{Sources}
\begin{itemize}
\item Local UnsolvedMath registry, permanent ID 30000203, exact source code catalogue entry; \emph{Minimal Subdegrees of Twisted-Wreath Permutation Groups}. Proposed year 2005, status open. The original explicitly restricts $P,H$ to Theorem 3 and refers to the bound of Theorem 1; their full statements are not supplied.
\item \emph{Groups and Geometries (2005)}, Oberwolfach Reports, DOI 10.4171/owr/2005/12, pp. 647--721, \url{https://doi.org/10.4171/owr/2005/12}.
\item Supplied literature triage checked 2026-08-21: ``original question'', \url{https://doi.org/10.4171/owr/2005/12}; ``primary minimal-subdegree study'', \url{https://doi.org/10.1515/9783110199741.75}. It identifies this as Question 4 and reports partial bounds and examples without a universal resolution. The dated OWR-derived curation requests independent verification of time-sensitive status and cited claims.
\end{itemize}
\end{document}
